% year: תשע"ז
% type: מבחן
% semester: א
% moed: ב
% number: 1
% points: 20
% categories: taylor
% source: exams/מבחנים/תשעז/מועד ב תשעז.pdf

\begin{question}
השתמשו בנוסחת טיילור כדי לחשב $\sin 167^\circ$ עם שגיאה לא גדולה מ-$10^{-5}$. הסבירו כיצד הערכתם את השגיאה.
\end{question}

\begin{hint}
$\sin167^\circ=\sin(180^\circ-13^\circ)=\sin13^\circ$. עברו לרדיאנים ופתחו את $\sin x$ סביב $0$.
\end{hint}

\begin{hint}
בפולינום מקלורן של $\sin x$ המקדם של $x^4$ הוא $0$, ולכן $P_3=P_4$ ואפשר להעריך את השארית בעזרת $R_4$.
\end{hint}

\begin{solution}
\textbf{רדוקציה.} $\sin167^\circ=\sin(180^\circ-13^\circ)=\sin13^\circ=\sin x_0$, כאשר
\[
x_0=\frac{13\pi}{180}\approx0.226893 \quad\text{(רדיאנים)}.
\]

\textbf{פולינום מקלורן.} עבור $f(x)=\sin x$: $f(0)=0$, $f'(0)=1$, $f''(0)=0$, $f'''(0)=-1$, $f^{(4)}(0)=0$, ולכן
\[
P_4(x)=P_3(x)=x-\frac{x^3}{6}.
\]

\textbf{הערכת השגיאה.} לפי משפט טיילור עם שארית לגרנז' מסדר $4$, קיים $c$ בין $0$ ל-$x_0$ כך ש-
\[
\sin x_0-P_4(x_0)=R_4(x_0)=\frac{f^{(5)}(c)}{5!}x_0^5=\frac{\cos c}{120}x_0^5 .
\]
מכיוון ש-$|\cos c|\le1$ ו-$0<x_0<0.23$:
\[
|R_4(x_0)|\le\frac{x_0^5}{120}<\frac{0.23^5}{120}\approx\frac{6.44\cdot10^{-4}}{120}\approx5.4\cdot10^{-6}<10^{-5}.
\]
(גם בהערכה גסה יותר $x_0<0.25$ מקבלים $\frac{0.25^5}{120}\approx8.1\cdot10^{-6}<10^{-5}$.)

\textbf{חישוב:}
\[
\sin167^\circ\approx x_0-\frac{x_0^3}{6}\approx0.226893-\frac{0.011680}{6}\approx0.226893-0.001947=0.224946.
\]

\textbf{תשובה:} $\sin167^\circ\approx\frac{13\pi}{180}-\frac16\left(\frac{13\pi}{180}\right)^3\approx0.22495$, עם שגיאה קטנה מ-$10^{-5}$ (הערך המדויק $0.2249511\ldots$).
\end{solution}
